% Chapter 1

\chapter{Εισαγωγή} % Main chapter title

\label{Chapter1} % For referencing the chapter elsewhere, use \ref{Chapter1} 

%----------------------------------------------------------------------------------------

% Define some commands to keep the formatting separated from the content 
\newcommand{\keyword}[1]{\textbf{#1}}
\newcommand{\tabhead}[1]{\textbf{#1}}
\newcommand{\code}[1]{\texttt{#1}}
\newcommand{\file}[1]{\texttt{\bfseries#1}}
\newcommand{\option}[1]{\texttt{\itshape#1}}

%----------------------------------------------------------------------------------------

\section{Το Πρόβλημα}

Η είσοδος στην καθημερινότητα μας, των συσκευών που συνδέονται στο Διαδύκτυο (IoT), έχει επηρεάσει αισθητά των όγκο των δεδομένων που συλλέγονται καθημερινά, καθώς και την ποικιλία τους. Το μεγαλύτερο ποσοστό αυτών των δεδομένων, είναι στην μορφή χρονοσειράς, δηλαδή μία ακολουθία από τιμές που έχουν καταγραφεί σε διαδοχικά χρονικά διαστήματα. Η ανίχνευση ανωμαλιών στις χρονοσειρές, είναι ένα πολύπλευρο πρόβλημα, το οποίο έχει κεντρίσει το ενδιαφέρον, τόσο των εταιριών αλλά και τον επιστημονικό κόσμο, καθώς σχετίζεται με υψηλότερης ποιότητας εφαρμογές στην υγεία, σε χρηματοοικονομικά συστήματα, σε περιβαλλοντολογικές παρακολουθήσεις, καθώς και σε πολλούς άλλους τομείς. Από την ανίχνευση ανωμαλιών μπορούμε να εξάγουμε πολύτιμες πληροφορίες, όπως π.χ. την ανίχνευση μία απάτης σε μία πιστωτική κάρτα, ή την ανίχνευση μίας σεισμικής δραστηριότητας, τον εντοπισμό ασυνήθιστης συμπεριφοράς ενός χρήστη σε ένα σύστημα κ.ο.κ

Τις τελευταίες δεκαετίες, έχουν προταθεί και αξιολογηθεί πληθόρα αλγορίθμων ανίχνευσης ανωμαλιών σε χρονοσειρές \cite{paparrizos2022tsb,gao2020ensemble,liu2008isolation,breunig2000lof,boniol2021sand,boniol2020automated,benecki2021detecting,boniol2021unsupervised,boniol2022theseus,boniol2021sand,yeh2016matrix,boniol2022series2graph,linardi2020matrix,yeh2018time,boniol2021sand_action}. Ωστόσο η επιλογή του κατάλληλου αλγορίθμου και η ρύθμιση των υπερπαραμέτρων του, είναι ένα πολυσύνθετο πρόβλημα, το οποίο εξαρτάται από το είδος της χρονοσειράς, των τύπο των ανωμαλιών, και απαιτεί πολύ χρόνο και υπολογιστικό κόστος.

Επίσης η έλλειψη ύπαρξης δεδομένων με ετικέτες για την εκπαίδευση των αλγορίθμων, μας οδηγεί σε μία μη-επιβλεπόμενη τεχνική για την ανίχνευση ανωμαλιών.

Τέλος, η έλλειψη ύπαρξης ενός καθολικού συστήματος σημείου αναφοράς (benchmark), το οποίο θα αξιολογεί τις διάφορες τεχνικές ανίχνευσης ανωμαλιών που έχουν ή πρόκειται να προταθούν, καθώς και μίας αμερόληπτης τεχνικής πέραν των παραδοσιακών, συνθέτουν ακόμα δύο προβλήματα προς εξερεύνηση.


%----------------------------------------------------------------------------------------

\section{Προτεινόμενες Λύσεις}

Τα τελευταία χρόνια (από το \date{2022} και μετά), έχουν προταθεί κάποια συστήματα, τα οποία έρχονται στοχεύουν στην επίλυση των προαναφερθέντων προβλημάτων.

\subsection{TSB-UAD Σύστημα σημείου αναφοράς (benchmark) για την ανίχνευση ανωμαλιών σε χρονοσειρές}

To 2022, προτάθηκε το \enquote{TSB: A Time Series Benchmark for Anomaly Detection} \cite{paparrizos2022tsb}, το οποίο αποτελείται από από μία συλλογή  αποτελούμενη από 13766 χρονοσειρές εκ των οποίων, οι 1980 προέρχονται από πραγματικά δεδομένα, οι 958 από τεχνητά δεδομένα, και 10828 συνθετικές χρονοσειρές κλιμακούμενης δυσκολίας.

Είναι ένα end-to-end benchmark, για μονοδιάστατες χρονοσειρές, το οποίο περιλαμβάνει μια ευρεία γκάμα τύπων ανωμαλιών. Προσπαθεί να παρέχει μία εννοιαία πλατφόρμα, σύγκρισης μεθόδων σε διαφορετικά ρεαλιστικά σενάρια, και να παρέχει ένα σύνολο από εργαλεία για την αξιολόγηση των αλγορίθμων ανίχνευσης ανωμαλιών σε χρονοσειρές.

\subsection{AutoTSAD: Ένα αυτοματοποιημένο σύστημα για την ανίχνευση ανωμαλιών σε χρονοσειρές}

To 2024, προτάθηκε το \enquote{AutoTSAD: Unsupervised Holistic Anomaly Detection for Time Series Data} \cite{schmidl2024autotsad} το οποίο είναι ένα μη-επιβλεπόμενο συνδιαστικό σύστημα ανίχνευσης ανωμαλιών σε μονοδιάστατες χρονοσειρές. Βασικά χαρακτηριστικά του έιναι ότι δεν απαιτεί την ύπαρξη ετικετών για την αξιολόγηση των αλγορίθμων, βασίζεται στον μηχανισμό Optuna \cite{akiba2019optuna} για την εύρεση των βέλτιστων υπερπαραμέτρων, και χρησιμοποιεί ένα ευρή φάσμα αλγορίθμων ανίχνευσης ανωμαλιών. Επίσης σε αντίθεση με άλλα συστήματα, το AutoTSAD επιστρέφει μία συνάθροιση βαθμολογιών, και όχι μία κατάταξη των αλγορίθμων. 

Τέλος μπορεί να λειτουργήσει και χωρίς την χρήση του μηχανισμού Optuna, παρά μόνο με προκαθορισμένες τιμές για τις υπερπαραμέτρους του κάθε αλγορίθμου, λόγω του τεράστιου απαιτούμενου χρόνου για την εύρεση των βέλτιστων υπερπαραμέτρων.

%----------------------------------------------------------------------------------------

\section{Σκοπός Πτυχιακής}
\label{purpose}

Σκοπός αυτής της πτυχιακής εργασίας, είναι η αξιολόγηση του AutoTSAD, το οποίο είναι ένα χρονικά απαιτητικό σύστημα, και να συγκριθέι η επίδοσή του, με μεμονωμένους αλγόριθμους, οι οποίοι έχουν καλή απόδοση ο καθένας σε διαφορετικά είδη τύπου ανωμαλιών, όπως ο MP \cite{linardi2020matrix}, ο IForest \cite{liu2008isolation}, ο POLY \cite{li1997poly} και NORMA \cite{mahin2024norma}, να παρουσιαστούν τα αποτελέσματα, και να αξιολογηθεί αν υπάρχει σοβαρή διαφορά η οποία να δικαιολογεί το χρονικό κόστος, πάντα με την χρήση αμερόληπτων τεχνικών, όπως η VUS-ROC\cite{paparrizos2022volume}.


%----------------------------------------------------------------------------------------

